\chapter{Expectation, Variance and Moments}\label{ch:expectation-moments}

\chapterinfo{STAT2003, STAT1301}{\cref{ch:probability-spaces,ch:random-variables}}

% Sources: STAT2003 handwritten notes (pp. 10-13, 63); STAT2003 course notes, Section 2.3;
% STAT1301 course notes (2026), Section 1.9.

All the probability information about a random variable is contained in its CDF, or in its pmf or pdf. Even so, it is often useful to consider numerical characteristics of a random variable. This chapter covers the most important ones --- the expectation, the variance and the moments --- and ends with two inequalities that bound probabilities in terms of these characteristics.

% ==========================================================================
\section{Expectation}\label{sec:expectation}
% ==========================================================================

The expectation of a random variable is a kind of ``weighted average'' of the values it can take.

\begin{definition}[Expectation]\label{def:expectation}
The \emph{expectation} (or \emph{expected value}, or colloquially the \emph{mean}) of a random variable $X$, written $\E(X)$ (also $\E[X]$ or $\E X$), is
\[
  \E(X) = \begin{cases}
    \displaystyle \sum_x x \, \Prob(X = x) = \sum_x x f(x) & \text{if $X$ is discrete with pmf $f$}, \\[10pt]
    \displaystyle \int_{-\infty}^{\infty} x f(x) \dd x & \text{if $X$ is continuous with pdf $f$}.
  \end{cases}
\]
The expectation of $X$ is sometimes written $\mu_X$. It is assumed that the sum or integral is well defined, and possibly infinite.
\end{definition}

\begin{remark}
The expectation is not the mean of a set of numbers, but the mean of a probability distribution.
\end{remark}

\begin{example}[Tossing a fair die]\label{ex:expectation-fair-die}
Find $\E(X)$ if $X$ is the outcome of a toss of a fair die.
\end{example}

\begin{solution}
$\Prob(X = 1) = \cdots = \Prob(X = 6) = \tfrac16$, so
\[
  \E(X) = 1 \cdot \tfrac16 + 2 \cdot \tfrac16 + \cdots + 6 \cdot \tfrac16 = \tfrac72.
\]
Note that $\E(X)$ need not be a possible outcome of the experiment.
\end{solution}

One way to interpret the expectation is as a long-run average payout. Suppose a game pays $X$ dollars, the largest face value of two dice, and costs $d$ dollars to play. Over many games, the fraction of games with maximum $1, 2, \ldots, 6$ is $\Prob(X = 1), \ldots, \Prob(X = 6)$, so the long-run average payout is the weighted sum (\cref{ex:max-two-dice})
\[
  \E(X) = 1 \times \tfrac{1}{36} + 2 \times \tfrac{3}{36} + 3 \times \tfrac{5}{36} + 4 \times \tfrac{7}{36} + 5 \times \tfrac{9}{36} + 6 \times \tfrac{11}{36} = \tfrac{161}{36} \approx 4.47.
\]
The game is fair if the long-run average profit $\E(X) - d$ is zero, so a fair fee is $d = \E(X)$ dollars.

Another interpretation of the expectation is as a centre of mass. Imagine point masses with weights $p_1, \ldots, p_n$ placed at positions $x_1, \ldots, x_n$ on the real line (\cref{fig:centre-of-mass}). The centre of mass, the place where we can ``balance'' the weights, is
\[
  x_1 p_1 + \cdots + x_n p_n,
\]
which is exactly the expectation of a discrete random variable taking the values $x_1, \ldots, x_n$ with probabilities $p_1, \ldots, p_n$ (compare to moments in physics). This interpretation gives the following.

\begin{figure}[H]
\centering
\begin{tikzpicture}
  \draw[thick] (0,0) -- (8,0);
  \foreach \x/\w/\l in {0.8/0.35/1, 2.6/0.5/2, 6.8/0.3/n} {
    \fill[thmcolor] (\x-\w/2,0) rectangle (\x+\w/2,\w*1.6);
    \node[font=\small] at (\x,\w*1.6+0.3) {$p_{\l}$};
    \node[font=\small] at (\x,-0.35) {$x_{\l}$};
  }
  \node at (4.6,0.3) {$\cdots$};
  \fill[remcolor] (3.3,0) -- (3.0,-0.5) -- (3.6,-0.5) -- cycle;
  \node[font=\small] at (3.3,-0.85) {$\E(X)$};
\end{tikzpicture}
\caption{The expectation as a centre of mass: weights $p_i$ at positions $x_i$ balance at $\E(X)$.}
\label{fig:centre-of-mass}
\end{figure}

\begin{proposition}[Expectation of a symmetric distribution]\label{prop:symmetric-expectation}
If the pmf or pdf of $X$ is symmetric about $c$, that is, $f(c + y) = f(c - y)$ for all $y$, then $\E(X) = c$ (provided the expectation is finite).
\end{proposition}

\begin{proof}
We give the discrete case:
\begin{align*}
  \E(X) &= c f(c) + \sum_{x > c} x f(x) + \sum_{x < c} x f(x)
        = c f(c) + \sum_{y > 0} (c + y) f(c + y) + \sum_{y > 0} (c - y) f(c - y) \\
       &= c f(c) + \sum_{y > 0} c f(c + y) + \sum_{y > 0} c f(c - y)
        = c \sum_x f(x) = c,
\end{align*}
where the third equality uses $f(c+y) = f(c-y)$ to cancel the terms $y f(c+y)$ and $-y f(c-y)$.
\end{proof}

\begin{example}[A continuous expectation]\label{ex:expectation-2x}
Let $X$ have pdf $f(x) = 2x$ for $x \in (0, 1)$, and $f(x) = 0$ otherwise. Then
\[
  \E(X) = \int_{-\infty}^{\infty} x f(x) \dd x = \int_0^1 2x^2 \dd x = \tfrac23 x^3 \Big|_0^1 = \tfrac23.
\]
\end{example}

% ==========================================================================
\section{Expectation of a function of a random variable}\label{sec:expectation-of-function}
% ==========================================================================

If $X$ is a random variable, then a function of $X$, such as $X^2$ or $\sin(X)$, is also a random variable. Its expected value is the weighted average of the values the function can take.

\begin{theorem}[Law of the unconscious statistician, LOTUS]\label{thm:lotus}
Let $g$ be a real-valued function. If $X$ is discrete with pmf $f$, then
\[
  \E(g(X)) = \sum_x g(x) f(x).
\]
If $X$ is continuous with pdf $f$, then
\[
  \E(g(X)) = \int_{-\infty}^{\infty} g(x) f(x) \dd x.
\]
\end{theorem}

\begin{proof}
We prove the discrete case only. Let $Y = g(X)$, which is again a random variable. Its pmf is
\[
  f_Y(y) = \Prob(Y = y) = \Prob(g(X) = y) = \sum_{x : g(x) = y} \Prob(X = x) = \sum_{x : g(x) = y} f_X(x).
\]
So the expectation of $Y$ is
\[
  \E(Y) = \sum_y y f_Y(y) = \sum_y y \sum_{x : g(x) = y} f_X(x) = \sum_y \sum_{x : g(x) = y} g(x) f_X(x) = \sum_x g(x) f_X(x). \qedhere
\]
\end{proof}

\begin{example}[Expected square of a fair die]\label{ex:expected-square-die}
For the outcome $X$ of a toss of a fair die,
\[
  \E(X^2) = 1^2 \cdot \tfrac16 + 2^2 \cdot \tfrac16 + 3^2 \cdot \tfrac16 + \cdots + 6^2 \cdot \tfrac16 = \tfrac{91}{6}.
\]
\end{example}

An important consequence of \cref{thm:lotus} is that expectation is linear.

\begin{theorem}[Linearity of expectation]\label{thm:linearity-single}
Let $h(x) = a_1 h_1(x) + a_2 h_2(x) + \cdots + a_n h_n(x)$, where $a_i \in \R$ and $h_i \colon \R \to \R$ for $i = 1, \ldots, n$. Then
\[
  \E(h(X)) = \sum_{i=1}^n a_i \E(h_i(X)).
\]
In particular, for any real numbers $a$ and $b$ and functions $g$ and $h$,
\begin{enumerate}
  \item $\E(aX + b) = a \E(X) + b$
  \item $\E\bigl(g(X) + h(X)\bigr) = \E(g(X)) + \E(h(X))$
\end{enumerate}
\end{theorem}

\begin{proof}
In the continuous case, by \cref{thm:lotus},
\begin{align*}
  \E(h(X)) &= \int_{-\infty}^{\infty} h(x) f(x) \dd x = \int_{-\infty}^{\infty} \sum_{i=1}^n a_i h_i(x) f(x) \dd x \\
  &= \sum_{i=1}^n a_i \int_{-\infty}^{\infty} h_i(x) f(x) \dd x = \sum_{i=1}^n a_i \E(h_i(X)).
\end{align*}
The discrete case is the same with sums in place of integrals. For example, statement 1 follows from
\[
  \E(aX + b) = \sum_x (ax + b) f(x) = a \sum_x x f(x) + b \sum_x f(x) = a \E(X) + b. \qedhere
\]
\end{proof}

% ==========================================================================
\section{Variance and moments}\label{sec:variance-moments}
% ==========================================================================

The variance of $X$ measures the spread, or dispersion, of its distribution.

\begin{definition}[Variance, standard deviation]\label{def:variance}
The \emph{variance} of a random variable $X$, written $\Var(X)$ or $\sigma_X^2$, is
\[
  \Var(X) = \E\bigl((X - \mu)^2\bigr), \qquad \text{where } \mu = \E(X).
\]
The \emph{standard deviation} of $X$ is $\sigma_X = \sqrt{\Var(X)}$.
\end{definition}

\begin{definition}[Moments]\label{def:moments}
The number $\E(X^r)$ is the \emph{$r$th moment} of $X$. With $\mu = \E(X)$, the number $\E\bigl((X - \mu)^k\bigr)$ is the \emph{$k$th central moment} of $X$.
\end{definition}

Here $\E(X)$ is the mean, and $\E(X^2)$ is related to the variance. The first central moments are
\begin{align*}
  \text{1st:} \quad & \E(X - \mu) = \E(X) - \mu = \mu - \mu = 0, \\
  \text{2nd:} \quad & \E\bigl((X - \mu)^2\bigr) = \Var(X).
\end{align*}
So the second central moment is the variance. The third and fourth central moments are related to the skewness and kurtosis of the distribution.

\begin{theorem}[Properties of the variance]\label{thm:variance-properties}
For any random variable $X$ and $a, b \in \R$,
\begin{enumerate}
  \item $\Var(X) = \E(X^2) - (\E X)^2$
  \item $\Var(aX + b) = a^2 \Var(X)$
\end{enumerate}
Consequently, $\sigma_{aX + b} = \abs{a}\,\sigma_X$.
\end{theorem}

\begin{proof}\leavevmode
\begin{enumerate}[leftmargin=*]
  \item[1.] Write $\E(X) = \mu$, so that $\Var(X) = \E\bigl((X - \mu)^2\bigr) = \E(X^2 - 2\mu X + \mu^2)$. By linearity of expectation (\cref{thm:linearity-single}), this equals $\E(X^2) - 2\mu \E(X) + \mu^2 = \E(X^2) - \mu^2$.
  \item[2.] The expectation of $aX + b$ is $a\mu + b$. So
\[
  \Var(aX + b) = \E\bigl((aX + b - (a\mu + b))^2\bigr) = \E\bigl(a^2 (X - \mu)^2\bigr) = a^2 \E\bigl((X - \mu)^2\bigr) = a^2 \Var(X).
\]
Taking square roots gives the statement about the standard deviation. \qedhere
\end{enumerate}
\end{proof}

% ==========================================================================
\section{Markov inequalities}\label{sec:markov-chebyshev}
% ==========================================================================

For an event $A$, let $\ind_A$ be the \emph{indicator} of $A$: the random variable that equals $1$ if $A$ occurs and $0$ otherwise. It takes the value 1 with probability $\Prob(A)$, so
\[
  \E(\ind_A) = \Prob(A).
\]

\begin{theorem}[Generalised Markov inequality]\label{thm:markov-general}
Let $h \colon \R \to [0, \infty)$. Then for every $a > 0$,
\[
  \Prob\bigl(h(X) \ge a\bigr) \le \frac{\E(h(X))}{a}.
\]
\end{theorem}

\begin{proof}
Fix $a > 0$ and let $A = \set{h(X) \ge a}$. On $A$, $h(X) \ge a = a \ind_A$; off $A$, $h(X) \ge 0 = a \ind_A$. So $h(X) \ge a \ind_A$ always, and therefore
\[
  \E(h(X)) \ge \E(a \ind_A) = a \Prob(A) = a \Prob\bigl(h(X) \ge a\bigr). \qedhere
\]
\end{proof}

\begin{corollary}[Markov inequality]\label{cor:markov}
For every $a > 0$,
\[
  \Prob(\abs{X} \ge a) \le \frac{\E(\abs{X})}{a}.
\]
In particular, if $X$ is non-negative, then $\Prob(X \ge a) \le \E(X) / a$.
\end{corollary}

\begin{proof}
Take $h(x) = \abs{x}$ in \cref{thm:markov-general}.
\end{proof}

\begin{corollary}[Chebyshev inequality]\label{cor:chebyshev}
Let $X$ have mean $\mu$. For every $a > 0$,
\[
  \Prob(\abs{X - \mu} \ge a) \le \frac{\Var(X)}{a^2}.
\]
\end{corollary}

\begin{proof}
Apply \cref{thm:markov-general} with $h(x) = (x - \mu)^2$, and $a^2$ in place of $a$:
\[
  \Prob(\abs{X - \mu} \ge a) = \Prob\bigl((X - \mu)^2 \ge a^2\bigr) \le \frac{\E\bigl((X - \mu)^2\bigr)}{a^2} = \frac{\Var(X)}{a^2}. \qedhere
\]
\end{proof}

\begin{remark}
For a non-negative continuous random variable with pdf $f$, the Markov inequality can also be seen directly:
\[
  \E(X) = \int_0^a x f(x) \dd x + \int_a^\infty x f(x) \dd x \ge \int_a^\infty x f(x) \dd x \ge \int_a^\infty a f(x) \dd x = a \Prob(X \ge a).
\]
\end{remark}
